\subsection{样本组成变化与分类型前沿}
总体经验前沿是按月份对当期样本分布取分位数，它并不对应一个固定模型随时间重复训练的轨迹。设第 $t$ 月不同模型类型的比例为 $\pi_{kt}$、各类型能力分布为 $H_{kt}(x)$，则总体分布为
\begin{equation}
H_t(x)=\sum_k\pi_{kt}H_{kt}(x),\qquad
F_t=\inf\{x:H_t(x)\geq0.9\}.
\label{eq:q4-mixture-frontier}
\end{equation}
一般情况下，总体第 $90$ 百分位不等于各类型第 $90$ 百分位的加权平均。即使每类能力分布保持不变，只要类型比例改变，总体前沿也可能移动；反过来，某一类型前沿上升，也未必使总体分位数等幅上升。因此，将所有月度变化直接解释为纯算法进步，会把样本构成效应混入趋势。

\begin{table}[!htbp]
\centering
\songti\zihao{-4}
\setlength{\tabcolsep}{4pt}
\caption{按模型类型分层的前沿样本与线性斜率}
\label{tab:q4-type-frontiers}
\begin{tabular}{lrrrr}
\toprule
类型 & 总记录 & 月最少 & 月最多 & 斜率（点/年） \\
\midrule
对话/微调 & 1589 & 80 & 234 & +18.17 \\
合并 & 674 & 19 & 116 & +27.04 \\
预训练 & 239 & 10 & 74 & -7.30 \\
\bottomrule
\end{tabular}
\end{table}


表~\ref{tab:q4-type-frontiers} 使用与主前沿相同的筛选和计分规则，分别统计三类主要模型。“其它”类仅有 $7$ 条记录，不为其拟合趋势。预训练组的月度样本量最少为 $10$，前沿波动也较大；该组的负拟合斜率不表示预训练技术随时间退步，而是当前观测窗口、当期提交对象和样本量共同决定的描述。对话或微调、合并模型占据绝大多数记录，其走势对总体前沿的影响更大。

\begin{figure}[!htbp]
\centering
\includegraphics[width=0.90\textwidth]{fig_q4_type_frontiers.pdf}
\caption{同一筛选规则下的分类型月度能力前沿；各曲线描述当月提交样本}
\label{fig:q4-type-frontiers}
\end{figure}

回归分析中的类型控制给出同样的提醒：加入类型虚拟变量后，年份系数由 $12.14$ 降至 $9.96$ 点/年，说明类型构成可以解释一部分时间关联，但仍不能据此把剩余年份系数认定为纯技术效应。模型家族、重复微调、参数来源、训练数据和评测版本都可能造成残余依赖。当前 HC3 标准误针对异方差调整，没有将同一家族的多条记录视为一个独立簇；其数值不应被解释为已经覆盖所有抽样相关性。

\subsection{滚动回测中的可用信息与预测失误}
前沿外推必须先通过只使用过去信息的检验。第 $j$ 次测试将前 $j-1$ 个有效月度点用于拟合，在第 $j$ 月预测前沿；不能先用全部月份估计斜率，再把历史点称为测试集。本文从 $5$ 个训练月份开始，依次得到 $5$ 个一步预测，并以最近一期真实前沿作为简单基线。表~\ref{tab:q4-backtest} 给出每次预测，而不仅报告平均误差。

\begin{table}[!htbp]
\centering
\songti\zihao{-4}
\setlength{\tabcolsep}{4pt}
\caption{扩展窗口的一步预测与实测前沿}
\label{tab:q4-backtest}
\begin{tabular}{rrrrr}
\toprule
测试时间 & 实测 & 趋势预测 & 末值预测 & 趋势误差 \\
\midrule
2024.875 & 34.592 & 37.039 & 34.670 & +2.447 \\
2024.958 & 39.953 & 37.518 & 34.592 & -2.435 \\
2025.042 & 39.986 & 40.670 & 39.953 & +0.684 \\
2025.125 & 40.779 & 42.350 & 39.986 & +1.571 \\
2025.208 & 41.887 & 43.616 & 40.779 & +1.729 \\
\bottomrule
\end{tabular}
\end{table}


均方根误差和平均绝对误差分别为
\begin{equation}
\operatorname{RMSE}=\sqrt{\frac{1}{m}\sum_{j=1}^{m}(\hat F_j-F_j)^2},
\qquad \operatorname{MAE}=\frac{1}{m}\sum_{j=1}^{m}|\hat F_j-F_j|.
\label{eq:q4-validation-errors}
\end{equation}
RMSE 对较大的预测偏差给予更高权重；平均有符号误差则检查是否持续高估或低估，不能替代绝对误差。当前趋势模型 MAE 为 $1.773$，RMSE 为 $1.889$，平均有符号误差为 $+0.799$。末值延续的 MAE 为 $1.475$。两者使用完全相同的测试月份，因而比较不受测试样本数量差异影响。

\begin{figure}[!htbp]
\centering
\includegraphics[width=0.88\textwidth]{fig_q4_rolling_validation.pdf}
\caption{扩展窗口的一步前沿预测与误差；训练窗口仅包含测试月份之前的数据}
\label{fig:q4-rolling-validation}
\end{figure}

图~\ref{fig:q4-rolling-validation} 将前沿的跳跃、平台与趋势预测放在同一时间轴上。一个较平稳的后续月份可能使末值基线表现良好，而前期上升趋势会继续向上外推；反之，出现跳跃时末值延续也可能明显低估。当前只有五次检验，无法为两类方法的长期优劣作稳定排序，但已经足以提醒：线性趋势在这组数据上并没有获得优于简单基线的短期预测证据。

因此，本文保留趋势情景的用途在于展示不同能力增长假设的后果，而不是宣称已经找到预测未来的最佳模型。若要把这一情景工具改为部署预测，应扩大历史窗口，按滚动起点评估多个预测期限，并在模型选择之后保留最终未参与调参的时间段。对区间预测还需检查实际覆盖率与区间宽度；只给出 Bootstrap 分位带，不能保证远期观测以相应概率落在其中。

\subsection{桥接与情景区间所覆盖的不确定性}
损失--能力桥接和能力前沿外推需要分别解释。对桥接，单调性约束保证在同一拟合层内损失减小不会降低预测能力，但这种方向来自函数约束，并不是额外的验证结果。高可比层的留一 RMSE 与均值基线之比约为 $1.10$，中可比层约为 $0.84$；这一无量纲比较表明，两层的预测价值并不一致。层间原始误差数值相差很大，也与能力值的变化范围不同有关，不能据此说高可比层的泛化能力远好于中可比层。

对前沿，在尚未裁剪到 $[0,100]$ 之前，令预测期限为 $u=\Delta/12$ 年，则抽样预测量的方差可展开为
\begin{align}
\operatorname{Var}(F_{s,\Delta})={}&\operatorname{Var}(\hat a_T)
+s^2u^2\operatorname{Var}(\hat\beta)
+2su\operatorname{Cov}(\hat a_T,\hat\beta)
+\operatorname{Var}(\epsilon),
\label{eq:q4-predictive-variance}
\end{align}
其中最后一项使用独立抽取的月度预测残差。联合重拟合保留了截距与斜率的协方差；若只扰动斜率而固定锚点，会遗漏前两者共同估计产生的变化。取经验分位数和边界裁剪之后，区间不再简单由正态方差决定，但该展开能说明各误差来源分别通过什么通道影响预测。

当 $s=0$ 时，式~\eqref{eq:q4-predictive-variance} 中含期限的项消失，故本实现的 $12$ 月与 $24$ 月停滞情景区间相同。这是“固定水平加同分布月度波动”的建模选择，不是越远期越没有额外风险。若未来存在随机游走、结构突变或评测改版，停滞情景的风险也可能随期限增加；当前区间并未包含这些机制。维持与放缓情景之间的差异属于假设差异，也不应并入单一情景的 Bootstrap 置信说明。

此外，C8 原始任务分数与榜单六维归一化分数存在口径差别。原始数学均值较低可用于描述该指标上的观测分布，却不能直接证明数学比所有其它任务更难；各任务的机会水平、评分上限和测试样本组成不同。使用桥接或前沿结果制定资源策略时，应先确定关心的是综合均分、某一任务的短板，还是某种部署约束下的表现，再选择对应的目标与验证集。
